\chapter{Implementation Comments}\label{sec:ImplementationComments}
Implementation comments should help to understand the code. This can be achieved in various ways.

\section{Comments to organize the Code}
When I started to learn programming, our teachers encouraged us to start with flow diagrams and alike, to sketch out the basic structure of the program that was refined in the later steps.

Later this sketching out was already done as \bdq{source code}, so that the first \bdq{version} of a program (or method\index{method}/function\index{function}) looked like this (here for the C\index{C} programming language):
\keeplisting{14}
\begin{lstlisting}[language=C]
int main( int argc, char *argv[] )
{
    //---* Parse the command line arguments *------------------------
    
    //---* Loop over the input file(s) *-----------------------------
    
    //---* Open the file *-------------------------------------------
    
    //---* Process each line *---------------------------------------
    
    //---* Merge the processing results *----------------------------
    
    //---* Write the output file *-----------------------------------
    
    //---* Do the housekeeping *-------------------------------------
}   //  main()
\end{lstlisting}

This compiles well, but does obviously nothing~…

As the code is further refined, certain steps are moved into separate methods or functions. Nevertheless, I leave the corresponding comment in the code, even though it is no longer strictly necessary. Among other things, it helps during development to keep track what is still needed.

%------------------------------------------------------------------------------

\section{Comments for Local Variables}\label{sec:LocalVariableComments}
The name for a variable should explain the function of that variable in the given context (refer to chapter \tqfullvref{sec:NamesForLocalVariables}), and for short sequences of code, that meaning can be easily deducted from the context. Consequently it should not be required to write comments for local variable (different from the parameters of a method – see \tqfullvref{sec:MethodParameterComments}). Nevertheless, sometimes it makes sense to provide that kind of additional information.

Usually, you will use trailing comments to provide the documentation for local variables:
\index{Gauss@Gauß, Carl Friedrich}\index{Lichtenberg, Heiner}\index{java.lang!IllegalArgumentException}\index{java.time!LocalDate}\index{java.time!LocalDate!minusDays()}\index{java.time!LocalDate!of()}\index{java.time!LocalDate!plusDays()}\index{java.time!Month!MARCH}\index{java.time!Year}\index{java.time!Year!getValue()}\index{java.time!Year!isAfter()}\index{java.time!Year!isBefore()}\index{org.tquadrat.foundation.lang!Objects}\index{org.tquadrat.foundation.lang!Objects!requireNonNullArgument()}
\begin{lstlisting}[numbers=left,caption={Gauß' Easter algorithm \autocite{WIKIPEDIA:DateOfEaster,WIKIPEDIA:Gaussche_Osterformel}},label={listing:GaussEaster}]
/**
 *  <p>{@summary This method calculates the date of Easter Sunday 
 *  for the given year.} The year has to be in the range from 1583 
 *  to 3900 (included).</p>
 *  <p>The resulting date is for the Gregorian calendar.</p>
 *  <p>The algorithm itself is not the original one published by Carl
 *  Friedrich Gauß first in 1800 (corrected version in 1816), but
 *  that one published by Heiner Lichtenberg in 1997.</p>
 *
 *  @thanks Carl Friedrich Gauß
 *  @thanks Heiner Lichtenberg  
 *
 *  @param  year    The year for which the date of Easter Sunday is 
 *      wanted for.
 *  @return The date of Easter Sunday in the given year.
 *  
 *  @see <a href="https://de.wikipedia.org/wiki/Gau%C3%9Fsche_Osterformel"> Gaußsche Osterformel</a>
 */
public static final LocalDate calcEasterDate( final Year year )
{
    /*
     * The explanation for the variables was taken from the German
     * Wikipedia article and translated by me. The original terms
     * are given in parenthesis.
     */
    final int x;  // The year
    final int k;  // The secular number (die Säkularzahl)   
    final int m;  // The secular moon shift (die säkulare Mondschaltung)
    final int s;  // The secular moon shift (die säkulare Sonnenschaltung)
    final int a;  // The moon parameter (der Mondparameter)
    final int d;  // The seed for the first full moon in spring (der Keim 
        // für den ersten Vollmond im Frühling)
    final int r:  // The calendar adjustment (die kalendarische
        // Korrekturgröße)
    final int og; // The Easter limit (die Ostergrenze)
    final int sz; // The first Sunday in March (der erste Sonntag im März)
    final int oe; // The distance of Easter Sunday from the Easter limit
        // – Easter distance in days (die Entfernung des Ostersonntags
        // von der Ostergrenze – Osterentfernung in Tagen)
    final int os; // The date of Easter Sunday as a March date with 
        // April 1st as March 32nd etc. (das Datum des Ostersonntags als 
        // Märzdatum – 32. März = 1. April usw.)
    
    //---* Check the arguments *-------------------------------------    
    if( requireNonNullArgument( year, "year" ).isBefore( Year.of( 1583 ) ) )
    {
        throw new IllegalArgumentException( "This method will work only for years greater than or equal to 1583" );
    }
    if( year.isAfter( Year.of( 3900 ) ) )
    {
        throw new IllegalArgumentException( "This method will work only for years less than or equal to 3900" );
    }

    //---* Lichtenberg's Easter formula *----------------------------
    x = year.getValue();
    k = x / 100;    
    m = 15 + (3 * k + 3) / 4 - (8 * k + 13) / 25;
    s = 2 - (3 * k + 3) % 4;
    a = x % 19;
    d = (19 * a + m) % 30;
    r = (d + a / 11) / 29;
    og = 21 + d - r;
    sz = 7 - (x + x / 4 + s) % 7;
    oe = 7 - (og - sz) % 7;
    os = og + oe;

    /*
     * Initialise the return value with the last day of February and 
     * add the calculated number of days.
     */
    final var retValue = LocalDate.of( x, MARCH, 1 )
        .minusDays( 1 )
        .plusDays( os );

    //---* Done *----------------------------------------------------
    return retValue;
}   //  calcEasterDate()
\end{lstlisting}

As you can see, for this use case it is acceptable to continue a long comment in the following line (see lines~31 and 32, for example) without using a block comment. 

But it is also acceptable here to use block comments\index{block comment} for the documentation of the local variables, as shown below:\keeplisting{13}
\begin{lstlisting}[numbers=left,firstnumber=30]
    …
    final int d;  /* The seed for the first full moon in spring (der Keim 
        für den ersten Vollmond im Frühling) */
    final int r:  /* The calendar adjustment (die kalendarische
        Korrekturgröße) */
    final int og; // The Easter limit (die Ostergrenze)
    final int sz; // The first Sunday in March (der erste Sonntag im März)
    final int oe; /* The distance of Easter Sunday from the Easter limit
        – Easter distance in days (die Entfernung des Ostersonntags von 
        der Ostergrenze – Osterentfernung in Tagen) */
    final int os; /* The date of Easter Sunday as a March date with 
        April 1st as March 32nd etc. (das Datum des Ostersonntags als 
        Märzdatum – 32. März = 1. April usw.) */
    …    
\end{lstlisting}
But in general, lengthy trailing comments like these should be avoided.

Providing a comment for local variables makes sense here because I used the original names from the algorithm as published by Heiner Lichtenberg\index{Lichtenberg, Heiner} in 1997, instead of finding my own descriptive names for them.

%------------------------------------------------------------------------------

\section{Comments for Method Arguments}\label{sec:MethodArgumentComments}
Chapter \tqvref{sec:MethodSignatures} discusses how the signature\index{signature} for a method\index{method} or a constructor\index{constructor} should be designed. The key requirements are:
\begin{itemize}[nosep]
	\item{Keep the parameter list short!}
	\item{Avoid (more than one) boolean parameters! See page \pageref{sec:BooleanArguments}.}
	\item{Avoid (too many) parameters of the same type! See page \pageref{sec:MultipleSameTypeArgs}.}
\end{itemize}

These requirements are often difficult – or sometimes even impossible – to meet. Furthermore, methods\index{method} or constructors\index{constructor} with excessively long or overly complex parameter lists are frequently found in 3\textsuperscript{rd} party libraries.

Coding the call to these methods\index{method} or constructors\index{constructor} are inherently error prone because mixing up two arguments or missing one can go undetected by the compiler.

In such cases, it makes perfect sense to add comments to the arguments. This will of course not help the compiler, but it documents the developer's intention and helps to fix the error.\footnote{Here it would help if \Java would support named arguments.}

\keeplisting{12}
Assume the following method declaration:
\begin{lstlisting}
/**
 *  <p>{@summary Draws a circle on the canvas.}</p>
 *
 *  @param  x   The {@code x} coordinate for the center of the circle.
 *  @param  y   The {@code y} coordinate for the center of the circle.
 *  @param  radius  The radius of the circle.
 */
public final void drawCircle( final double x, final double y, final radius )
{
    …
}
\end{lstlisting}

\keeplisting{5}
When this method\index{method} is called with variables, the meaning of the arguments is reasonably clear - or perhaps not\footnote{Keep in mind that it is not required that the variables are named as \lstinline|x|, \lstinline|y| and \lstinline|d|.}:
\begin{lstlisting}
drawCircle(
    x, y, // The center of the circle.
    d / 2 // The radius as half of the diameter.
);
\end{lstlisting}

\keeplisting{5}
But when called with literal values, a comment for each argument is strongly recommended:
\begin{lstlisting}
drawCircle(
    5.678, // The x coordinate for the center of the circle. 
    7.322, // The y coordinate for the center of the circle.
    12.44  // The radius for the circle.
);
\end{lstlisting}

The chapter \tqfullref{sec:MethodSignatures} provides several strategies to avoid this situation.

%------------------------------------------------------------------------------

%------------------------------------------------------------------------------






 
But if the code is suThese requirements are often difficult—or sometimes even impossible—to meet. Furthermore, methods with excessively long or overly complex parameter lists are frequently found in third-party libraries.
In such cases, it makes perfect sense to add comments to the arguments.pposed to be self-explanatory, why do I still need comments? Short answer: because not everyone has the same understanding of the programming language! At the two extremes, you have rookies on one side and experts on the other – with a wide range of levels in between. Not to mention your own future self.

Furthermore, comments should explain why the code is written the way it is rather than in another way. There is often more than one way to implement something, and a comment should clarify why that specific approach was chosen. This is particularly important when the obvious solution was not selected or when the chosen implementation deviates from established best practices.

%------------------------------------------------------------------------------

\chapter{Comments for Local Variables}\label{sec:LocalVariableComments}
\lipsum[1]









 – with the result, that newbies do not have any help to understand the code written by the experts.

Obviously, the advice is the other extrema – and equally bad than the first advice, mentioned above!

Consequently it means that writing proper comments remains a complex art, but it follows some rules, and for the rest, I will try to give some advice.

First, comments have to support the readability of source code. They can do this by structuring the source code into sections and by providing hints on how to interpret the code.

Then they should be used to give an overview on the code and to provide additional information that is not readily available in the code itself. But it has to be relevant for the current context. For example, information about how the corresponding package is built or in what directory it resides should not be included as a comment to a class, but it might be useful in the comment for that package.

The discussion of non-trivial or non-obvious design decisions is appropriate, but the duplication of information that is already present in (and clear from) the code must be avoided. It is too easy for redundant comments to get out of date when the code changes.

And finally, comments should be in full sentence and using a clear language. They have to be in English language\footnote{Of course you can decide for any other language as well, but we already discussed in \tqfullvref{sec:Language} why English is preferred}, and they should be grammatically correct and without typos.\footnote{…~but it is still much more important that there is at least \textit{some} comment than a correctly spelled one.} This improves readability and helps to avoid misunderstandings.

But the final answer regarding the audience is still missing; I will discuss this further in \tqfullvref{sec:CommentsWhen}.

\section{Implementation Comment Formats}
A Java source code file can have three \bsq{styles} of implementation comments that are described in detail in the following chapters:
\begin{itemize}[nosep]
\item{block comments}
\item{single-line comments}
\item{trailing or end-of-line comments}
\end{itemize}


\subsection{Block Comments}\label{sec:BlockComments}
Block comments are used to provide detailed descriptions of files, methods, data structures and algorithms – meaning that the text of the comment is longer than just one single line. Block comments may be used at the beginning of each block after the opening brace. They can also be used in other places, such as within methods. Block comments inside a function or method should be indented to the same level as the code they describe.

A block comment should be preceded by a blank line to set it apart from the rest of the code. If the block comment does not directly refer to the code line immediately after it, it should be followed by another blank line.

Next, the first line of the comment block has to remain empty, and the closing of the comment block has to be placed on a line of its own.

Some samples:
\begin{lstlisting}
{
    /*
     * Here is a sample of a block comment. Block comments are used
     * to provide detailed information about code internals.
     */
    Result value = retrieveResult( parameter );
    …

    /*
     * Here is another sample of block comment, somewhere in the
     * middle of a code block. Please note the blank line above!
     */
    …
}

// AVOID!
/*
 * This block comment is outside the code block it refers to. Block 
 * comments should be placed after the opening curly brace of the
 * block.
 */
{
    Result value = retrieveResult( parameter );
    …
 
/* The first line of the block comment should be left empty and the
 * comment should be indented in the same way as the code in the
 * block.
 */
    processResult( value, parameter );
    /* 
     * Here the empty line above the comment is missing …
     * … and the closing tag should be on a line of its own. */
\end{lstlisting}
 
%------------------------------------------------------------------------------
 
\subsection{Single-Line Comments}\label{sec:SingleLineComments}
Short comments can appear on a single line; they will be indented also to the level of the code that follows. Usually it should be written in the form of a headline:
\begin{lstlisting}
//---* Handle the condition *----------------------------------------
\end{lstlisting}
with the dashes ending on column~80.\footnote{A quick reminder: the sample code in this document uses a line length of 70, so the dashes in the comment line above ends an column~70.}

If a comment cannot be written in a single line, it should have the block comment format (see chapter \tqref{sec:BlockComments}). A single-line comment should be separated from the preceding code by a blank line. Only when the preceding line contains only the opening curly brace as in the \lstinline|if-then-else| sample below (lines~3 and 8), this blank line should be omitted. 

Here are some examples of single-line comments in Java code:
\begin{lstlisting}[numbers=left]
if( cache.contains( key ) )
{
    //---* Take the data from the cache *----------------------------
    …
}
else
{
    //---* Load the data from its original source *------------------
    …
}
…
ResultData resultData = executeService();

//---* Format the output for the UI *--------------------------------
formatResult( resultData );
…
\end{lstlisting}

Another form of the single line comment is the empty block comment:
\begin{lstlisting}
public interface Marker
{ /* No methods */ }
//  interface Marker

public class Extension extends Base
{ /* No implementation */ }
//  class Extension

private Constructor() { /* Just exists */ }

public void adapterMethod() { /* Does nothing */ }

public final void method()
{
    …

    try
    {
        …
    }
    catch( final MyException e ) { /* Exception deliberately swallowed /* }

    …
}   // method()
\end{lstlisting}

And finally, there are the “class termination comments” that repeats the name of the class after the closing curly brace of the class definition, as you may have noticed already in several other examples:
\begin{lstlisting}
public final class MyClass
{
    …
}
//  class MyClass

public final interface MyInterface
{
    public record InnerRecord( final int number )
    {
        …
    }
    //  record InnerRecord
    
    …
}
//  interface MyInterface
\end{lstlisting}

%------------------------------------------------------------------------------

\subsection{Trailing or End-Of-Line Comments}\label{sec:TrailingOrEndOfLineComments}
Very short comments can appear on the same line as the code they describe, but should be shifted right far enough to separate them from the statements – and then that comment should still end at or before column~80. If more than one short comment appears in a chunk of code, they should all be indented to the same tab setting.

Here's an example of a trailing comment in Java source code:\footnote{Of course that can be written easier as “\lstinline!final var retValue = (a == 2) || isPrime( a );!”, but then it would be difficult to place the comments as shown. Again, another form of comment would be still possible …}
\begin{lstlisting}
final boolean retValue;
if( a == 2 )
{
    retValue = true;            /* special case */
}
else
{
    retValue = isPrime( a );    /* never true for even a */
}
\end{lstlisting}

But it is more common to use “\verb#//#” instead of “\verb#/*…*/#” for these trailing comments:
\begin{lstlisting}
final boolean retValue;
if( a == 2 )
{
    retValue = true;            // special case
}
else
{
    retValue = isPrime( a );    // never true for even a
}
\end{lstlisting}

You will use trailing comments to provide the documentation for local variables\footnote{Usually, the name of that local variable should be sufficient (refer to chapter \tqfullvref{sec:NamesForLocalVariables}), or the meaning of that variable is obvious from the context, but sometimes it still make sense to provide that kind of additional information.}:

\index{Gauss@Gauß, Carl Friedrich}\index{Lichtenberg, Heiner}\index{java.lang!IllegalArgumentException}\index{java.time!LocalDate}\index{java.time!LocalDate!minusDays()}\index{java.time!LocalDate!of()}\index{java.time!LocalDate!plusDays()}\index{java.time!Month!MARCH}\index{java.time!Year}\index{java.time!Year!getValue()}\index{java.time!Year!isAfter()}\index{java.time!Year!isBefore()}\index{org.tquadrat.foundation.lang!Objects}\index{org.tquadrat.foundation.lang!Objects!requireNonNullArgument()}
\begin{lstlisting}[numbers=left,caption={Gauss' Easter algorithm\autocite{WIKIPEDIA:DateOfEaster,WIKIPEDIA:Gaussche_Osterformel}},label={listing:GaussEaster}]
/**
 *  <p>{@summary This method calculates the date of Easter Sunday 
 *  for the given year.} The year has to be in the range from 1583 
 *  to 3900 (included).</p>
 *  <p>The resulting date is for the Gregorian calendar.</p>
 *  <p>The algorithm itself is not the original one published by Carl
 *  Friedrich Gauß first in 1800 (corrected version in 1816), but
 *  that one published by Heiner Lichtenberg in 1997.
 *
 *  @thanks Carl Friedrich Gauß
 *  @thanks Heiner Lichtenberg  
 *
 *  @param  year    The year for which the date of Easter Sunday is 
 *      wanted for.
 *  @return The date of Easter Sunday in the given year.
 *  
 *  @see <a href="https://de.wikipedia.org/wiki/Gau%C3%9Fsche_Osterformel">Gaußsche Osterformel</a>
 */
public static final LocalDate calcEasterDate( final Year year )
{
    /*
     * The explanation for the variables was taken from the German
     * Wikipedia article and translated by me. The original terms
     * are given in parenthesis.
     */
    final int x;  // The year
    final int k;  // The secular number (die Säkularzahl)   
    final int m;  // The secular moon shift (die säkulare Mondschaltung)
    final int s;  // The secular moon shift (die säkulare Sonnenschaltung)
    final int a;  // The moon parameter (der Mondparameter)
    final int d;  // The seed for the first full moon in spring (der Keim 
        // für den ersten Vollmond im Frühling)
    final int r:  // The calendar adjustment (die kalendarische
        // Korrekturgröße)
    final int og; // The Easter limit (die Ostergrenze)
    final int sz; // The first Sunday in March (der erste Sonntag im März)
    final int oe; // The distance of Easter Sunday from the Easter limit
        // – Easter distance in days (die Entfernung des Ostersonntags
        // von der Ostergrenze – Osterentfernung in Tagen)
    final int os; // The date of Easter Sunday as a March date with 
        // March 32 as April 1 etc. (das Datum des Ostersonntags als 
        // Märzdatum – 32. März = 1. April usw.)
    
    //---* Check the arguments *-------------------------------------    
    if( requireNonNullArgument( year, "year" ).isBefore( Year.of( 1583 ) ) )
    {
        throw new IllegalArgumentException( "This method will work only for years greater than or equal to 1583" );
    }
    if( year.isAfter( Year.of( 3900 ) ) )
    {
        throw new IllegalArgumentException( "This method will work only for years less than or equal to 3900" );
    }

    //---* Lichtenberg's Easter formula *----------------------------
    x = year.getValue();
    k = x / 100;    
    m = 15 + (3 * k + 3) / 4 - (8 * k + 13) / 25;
    s = 2 - (3 * k + 3) % 4;
    a = x % 19;
    d = (19 * a + m) % 30;
    r = (d + a / 11) / 29;
    og = 21 + d - r;
    sz = 7 - (x + x / 4 + s) % 7;
    oe = 7 - (og - sz) % 7;
    os = og + oe;

    /*
     * Initialise the return value with the last day of February and 
     * add the calculated number of days.
     */
    final var retValue = LocalDate.of( x, MARCH, 1 )
        .minusDays( 1 )
        .plusDays( os );

    //---* Done *----------------------------------------------------
    return retValue;
}   //  calcEasterDate()
\end{lstlisting}

As you can see, for this use case it is acceptable to continue a long comment in the following line (see lines~31 and 32, for example) without using a block comment. But writing
\begin{lstlisting}[numbers=left,firstnumber=30]
    …
    final int d;  /* The seed for the first full moon in spring (der Keim 
        für den ersten Vollmond im Frühling) */
    final int r:  /* The calendar adjustment (die kalendarische
        Korrekturgröße) */
    final int og; // The Easter limit (die Ostergrenze)
    final int sz; // The first Sunday in March (der erste Sonntag im März)
    final int oe; /* The distance of Easter Sunday from the Easter limit
        – Easter distance in days (die Entfernung des Ostersonntags von 
        der Ostergrenze – Osterentfernung in Tagen) */
    final int os; /* The date of Easter Sunday as a March date with 
        March 32 as April 1 etc. (das Datum des Ostersonntags als 
        Märzdatum – 32. März = 1. April usw.) */
    …    
\end{lstlisting}
can be considered also as valid. But in general, lengthy trailing comments like these should be avoided.

Another use case for these kind of comments is to provide information about the arguments of a method call; usually, you should avoid method signatures\index{signature} where this is required (the formal parameter of the method should be sufficient to explain the argument), but sometimes it makes still sense:
\begin{lstlisting}
…
drawCircle(
    x, y, // The center of the circle
    d/2.0 // The radius of the circle, calculated from the diameter
);
…
\end{lstlisting}
Again, you should avoid lengthy comments, but if really required, you can break the lines as shown above.

The end comment for a method is also of this type; I used lots of these in the sample code already:
\begin{lstlisting}
public final void method()
{
    …
} //  method()
\end{lstlisting}

Also a long code block\footnote{But if it seems really necessary to add such a comment to a linear code block, you should consider to re-organise your code.} can be commented like this:
\begin{lstlisting}
{
    //---* Calculate the result *------------------------------------
    // Lots of code comes here ...
    …
} //  End of result calculation
\end{lstlisting}

But if the comment should only mark begin and end of the code block, without providing further information, you should consider to use a label, as explained below.

The end comments for the bodies of \lstinline|for|, \lstinline|while|, \lstinline|switch| and even \lstinline|if-then-else| blocks are a special case, as those blocks should be introduced by a label, and this label should be repeated as the end comment:
\index{java.lang!IllegalArgumentException}
\begin{lstlisting}
ScanLoop: for( final var s : lines )
{
    // Lots of code comes here ...
    …
}   //  ScanLoop:

ForeverLoop: while( true )
{
    // Lots of code comes here ...
    …
}   //  ForeverLoop:

TypeSwitch: switch( type )
{
    case TYPE_1 -> …
    …
    case TYPE_n -> …
    default -> throw new IllegalArgumentException()
}   //  TypeSwitch:

SpecialCaseTurnout: if( isSpecialCase() )
    // Lots of code comes here ...
    …
}   //  SpecialCaseTurnout:
else
{
    // Not so much code here
    …
}    
\end{lstlisting}
Note the colon at the end of each of the comments!

Do not use the introducing code line for the end comment, although you can find this pattern in many books and even in real world code! That introducing code line may change at some point in time and then the end comment does not have a corresponding starting line anymore.\footnote{This also means that you should be careful with removing or changing labels in the code.}

Avoid this:
\index{java.lang!IllegalArgumentException}
\begin{lstlisting}
// AVOID!!!
for( final var s : lines )
{
    // Lots of code comes here ...
    …
}   //  for( final var s : lines )

while( true )
{
    // Lots of code comes here ...
    …
}   //  while( true )

switch( type )
{
    case TYPE_1 -> …
    …
    case TYPE_n -> …
    default -> throw new IllegalArgumentException()
}   //  switch( type )

if( isSpecialCase() )
    // Lots of code comes here ...
    …
}   //  if( isSpecialCase() )
else
{
    // Not so much code here
    …
}    
\end{lstlisting}

%------------------------------------------------------------------------------

\section{Comments when?}\label{sec:CommentsWhen}
To repeat what was already said at the beginning of this chapter \tqfullref{sec:WritingProperComments}: The frequency of comments reflects the quality of the source code, but it can be “under-commented” or “over-commented”. This means that there is a “frequency band” for comments that has to be hit for \textit{good} quality code.

The advice “When you feel compelled to add a comment, consider rewriting the code to make it clearer” raises the question for whom the code has to made clearer, because the experts and the rookies are completely distinct audiences, with different needs in regard of comments in the source code.

Same with that other advice (“Even if you think, a comment might not be necessary, add it nevertheless”): it is the other extrema, and equally bad.

That leads us (again) to the conclusion, that writing good comments to source code is a complex art.

This means that this document cannot give you a complete rule set here, just some guidance.

That documentation comments are required everywhere possible, and how to write them was already covered in the chapter \tqfullvref{sec:DocumentationComments}, thus we focus here on the implementation comments.

\begin{enumerate}[label=P\arabic*.]
\item{When something looks like a mistake, write a comment.

This could be an empty code block, an exception that is not handled, a fall-through in \lstinline|switch| branch, or an assignment where a comparison is expected. Here the comment should make it clear that this “mistake” is also intentional.}

\item{When you break a rule, write a comment.

The \ngref{} say in bullet points~\ref{lst:ZoP:SpecialCases} and \ref{lst:ZoP:Practicality}: “\textit{Special cases aren't special enough to break the rules although practicality beats purity}”. You will encounter lots of occasions in your daily work where this seems to fit – according to your opinion! Others may think differently about that, but a comment at least tells them why you wrote the codes as you did.

And by the way: “practicality” does not mean “laziness in typing".}

\item{Don't let them guess, tell them the solution.

In bullet point~\ref{lst:ZoP:DontGuess} of the \ngref{} you are advised not to guess – that's an advice also for the readers of the code. Add a comment when you think that something can be interpreted in different ways.}

\item{Describe the short-cuts.

Assume the following case: you have to parse a string with some input data, and you know from the specification, that this input data has a determined content when the string a specified length, and that this length will never be possible for other contents.

In that case, you can skip the parse process after you determined the length of the string and just set the result. But you should add some lengthy comment to your code that explains that ‘magic’.

And you will have this kind of ‘short-cut’ also in less esoteric cases: in chapter \tqfullvref{sec:CheckingMethodParametersAndReturnValues}, I request that the arguments to methods needs to be validated before they are used. But when the method is \lstinline|private|, all the callers to that method are well known and under your control. Therefore you may have checked the arguments already on the caller and can omit the check in the method itself. Or you delegated the validation to another method – in both cases, add a comment.}

\end{enumerate}

%------------------------------------------------------------------------------

\subsection{Empty Code Blocks are forbidden}
That means that the sequence “\verb#{}#” may not occur anywhere in your source code. If there is really no code in a block, there have to be at least a comment, just to show that this is intentional. For a \lstinline|catch| block, there should be a good (at least some) explanation, why it is empty (see also chapter \tqfullref{sec:GeneralExceptionHandling}).

Some samples:
\begin{lstlisting}
//---* Skip input stream until then end *----------------------------
while( input.read() != EOF ) { /* empty */ }

// Better
//---* Skip input stream until then end *----------------------------
while( input.read() != EOF );
\end{lstlisting}
\begin{lstlisting}
//---* Spent CPU cycles on counting *--------------------------------
for( var i = 0; i < maxValue; ++i ) { /* empty */ }

// Better
//---* Spent CPU cycles on counting *--------------------------------
for( var i = 0; i < maxValue; ++i );
\end{lstlisting}
In this case it is likely that the compiler will optimise the loop, just be removing it, because the code does nothing useful.
\begin{lstlisting}
InputStream input = null; // null if file cannot be opened 
try
{
    input = new FileInputStream( "myFile" );
}
catch( final FileNotFoundException ignored ) { /* Deliberately ignored */ }

// Better
final InputStream input; // null if file cannot be opened
try
{
    input = new FileInputStream( "myFile" );
}
catch( final FileNotFoundException ignored ) 
{ 
    input = null;
}
\end{lstlisting}
\begin{lstlisting}
public MyClass() { /* Just exists */ }
\end{lstlisting}
\begin{lstlisting}
public interface MyInterface { /* No methods */ }
\end{lstlisting}
\begin{lstlisting}
@MountPoint
public Data customAction( final Data data ) { /* Does nothing */ }
\end{lstlisting}

The usage of the annotation \lstinline|@MountPoint|\index{mountpoint} \autocite{TQUADRAT_ORG_FOUNDATION_MOUNTPOINT} is described in the chapters \tqfullvref{sec:NonFinalClasses} and \tqfullvref{sec:NonFinalMethods}, a sample implementation can be found in chapter \tqvref{sec:MountPoint}.

%------------------------------------------------------------------------------

\subsection{Description of the used Algorithm}
A short description of the algorithm that was implemented by the current source code is always helpful; even just mentioning its name can make a difference when hunting a bug. But in case the implementation changes, such a comment needs to be adjusted as well.

The same is valid for the pattern that you implement with your code, although this belongs more often into the documentation comments.

%------------------------------------------------------------------------------

\subsection{Broken Rules}
Whenever you hurt or ignore one of the rules, guidelines, or recommendations from this document, it is worth a comment; same when you use a short-cut.

Some examples:
\begin{lstlisting}
value = 7.0 / a; // a != 0.0 was already checked above

o.execute(); // o != null was checked in foo.bar( o )

/*
 * bar.getQ() will never return null, so an explicit check on
 * null for q was omitted
 */
q = bar.getQ();
q.execute();
\end{lstlisting}
In all these samples, arguments or values are not checked for \lstinline|null| or zero, \textit{because this was already done elsewhere} – and not necessarily just one line above the current location. The comment tells a maintenance engineer that the check was made – at least in the initial version of the code – and that the root cause for the problem might be elsewhere (perhaps because someone eliminated that other location or at least the value check there).

Non-obvious class casts can be seen as a short-cut or a broken rule, but they should be always explained with a comment, in order to document that you knew what you did:
\begin{lstlisting}
private final void valueProcessor( List<Object> values )
{
    …
    Value value;
    for( Object o : values )
    {
        /*
         * We know that values can only contain Value objects;
         * otherwise this method would not have been called to
         * process the list.
         */
        value = (Value) o;
        …
}
\end{lstlisting}
This makes sense here because a check like
\begin{lstlisting}
if( o instanceof Value ) value = (Value) o;
\end{lstlisting}
is relatively expensive – especially if the contract for this \lstinline|private| method is that it is called only with lists containing \lstinline|Value| objects so that the check would not be positive only in very, very, very rare cases. And finally: what else can be done in cases where the object is not of the right type than throwing a ClassCastException? That's the same that is done by the code above in such a case, too.

First, the version
\begin{lstlisting}
if( o instanceof Value value )
{
    …
}
\end{lstlisting}
might be clearer, but it is not less expensive.

Second, having a list of \lstinline|Value| objects declared as a \lstinline|List<Object>| is bad design anyway – or legacy code (and still bad design).

A fall-through in a \lstinline|switch| is another broken rule and needs a comment, as said already in chapter \tqfullvref{sec:SwitchStatements}.

%------------------------------------------------------------------------------

\subsection{Swallowed Exceptions}
The \ngref{} says in bullet point~\ref{lst:ZoP:ErrorsMayNeverBeSwallowed} that errors may not be ignored and that exceptions may not be swallowed; this is most often done by an empty \lstinline|catch| block.

But people forget that \lstinline|catch| blocks like those below will also swallow the exception, although it does not look like that:
\begin{lstlisting}[numbers=left]
…
catch{ final Exception e }
{
    e.printStackTrace();
}    

…
catch{ final IOException e }
{
    m_Logger.atError()
        .withThrowable( e )
        .log( "Opening the file failed" );
}

…
catch( final NumberFormatException e )
{
    m_LastError = e;
}       
\end{lstlisting}
All this might be intended, but at a first glance it looks like a mistake. This mean that a comment is mandatory each time an exception is swallowed, no matter if silently or not, meaning the exception is logged somewhere. If the program is continued after an exception was caught, write a comment (see also chapter \tqfullref{sec:GeneralExceptionHandling}).

Also when you catch \lstinline|java.lang.Error| or a child class of it, or \lstinline|java.lang.Throwable|\index{java.lang!Throwable}, you should provide a comment why you are doing that (and why you think it is allowed in your special case~…).

%------------------------------------------------------------------------------

\subsection{Suppressed Warnings}
With the \lstinline|@SuppressWarnings| annotation\autocite{ORACLE_DOC_SUPPRESSWARNINGS_ANNOTATION}, you can switch off a compiler warning or error caused by code that is somehow “hurting the rules”. Basically, the \lstinline|@SuppressWarnings| annotation is a replacement for a comment about that deviation from the rules, an additional comment is required only in cases it is not obvious why that annotation was applied. More details on this are provided in \tqfullvref{sec:CompilerWarningsAndErrors}.

%------------------------------------------------------------------------------


%==============================================================================
% End of file!!
%==============================================================================

